\documentclass[11pt]{article}
\usepackage[margin=1in]{geometry}
\usepackage{amsmath,amssymb,booktabs,longtable,array}
\usepackage{xcolor}
\usepackage[colorlinks,linkcolor=blue!45!black,urlcolor=blue!45!black]{hyperref}

\newcommand{\F}{\textsuperscript{F}}
\newcommand{\Sec}{\textsuperscript{S}}
\newcommand{\T}{\mathsf{T}}
\newcommand{\R}{\mathbb{R}}

\setlength{\parindent}{0pt}
\setlength{\parskip}{5pt}

\begin{document}

References of the form ``Prop.~4.6'' point to the derivation document. Provenance of each
evidence item is marked \F{} if the full text was retrieved and the quoted statement read
directly, and \Sec{} if it rests on a secondary summary and has not been independently verified.
Every \Sec{} item must be confirmed before it is relied on in a submission.

\section{Summary}

\begin{longtable}{@{}clll@{}}
\toprule
\# & Result & Closest prior & Verdict\\
\midrule
\endhead
1 & Prop.~4.6, marginal field is linear & Hertrich et al., Thm.~3(i) & \textbf{Same}\\
2 & Lem.~4.4, the identity & Tsimpos et al., Lem.~24, Rmk.~2 & \textbf{Partial}\\
3 & Prop.~4.9, marginal matching, any $N$ & De Bortoli et al.; Wald--Steidl & \textbf{Partial}\\
4 & Thm.~4.11, pencil criterion & none & \textbf{No prior found}\\
5 & Cor.~4.13, commuting $\Rightarrow$ Brenier & Hertrich et al., Thm.~3(ii) & \textbf{Same}\\
6 & Cor.~4.21, Rmk.~4.22, degenerate source & Tsimpos et al., Prop.~10, Rmk.~8 & \textbf{Partial}\\
7 & Cor.~5.5, azimuth conservation & standard symmetry reduction & \textbf{Partial}\\
8 & Thm.~5.9, spherical landing law & none & \textbf{No prior found}\\
9 & Prop.~6.3, $\beta^\star$ exists iff $A<\cos j$ & none & \textbf{No prior found}\\
10 & \S7, the model $\mathcal M_2$ & De Bortoli et al., Eq.~(10) & \textbf{Same construction}\\
11 & Prop.~7.3, Rmk.~7.5, degeneracy at $t=0$ & Wang et al., Thms.~1--2 & \textbf{Partial}\\
12 & Meas.~7.6, $\mathcal M_1$--$\mathcal M_2$ inversion & none & \textbf{No prior found}\\
13 & Cor.~3.3, cosine rewards means & Blau \& Michaeli 2018 & \textbf{Same}\\
14 & Stability law refining Thm.~4.11 & none & \textbf{No prior found}\\
15 & Super-resolution as an instance & Delbracio \& Milanfar; Blau \& Michaeli & \textbf{Partial}\\
16 & Prop.~4.33, loss floor does not determine $M$ & Zhang et al., App.~K & \textbf{Partial}\\
\bottomrule
\end{longtable}

\section{Detail}

\subsection*{1. Prop.~4.6 --- Same}

\textbf{Prior.} J.~Hertrich, A.~Chambolle, J.~Delon, \emph{On the Relation between Rectified Flows
and Optimal Transport}, arXiv:2505.19712, Theorem~3(i).\F

\textbf{Statement quoted.} ``The minimizer $v_t=\arg\min_{w_t}\mathcal L(w_t|X_0,X_1)$ of the loss
function (1) is given by
$v_t(x)=\frac{1}{1-t}\big(((1-t)\Sigma_{01}+t\Sigma_1)\Sigma_t^{-1}-\mathrm{Id}\big)x$, where
$\Sigma_t=\mathrm{Cov}(X_t)=(1-t)^2\Sigma_0+(1-t)t(\Sigma_{01}+\Sigma_{10})+t^2\Sigma_1$.''
No independence is assumed; the formula holds for arbitrary cross-covariance.

\textbf{Identification.} Expanding $(1-t)\Sigma_{01}+t\Sigma_1-\Sigma_t$ and factoring $(1-t)$,
\[
v_t=\big[(1-t)(\Sigma_{01}-\Sigma_0)+t(\Sigma_1-\Sigma_{10})\big]\Sigma_t^{-1},
\]
which is $K(t)=G(t)V(t)^{-1}$ of Prop.~4.6 under their block convention
$\Sigma_{01}=C^{\T}$, $\Sigma_{10}=C$, with $\Sigma_t=V(t)$ identically.

\textbf{Action.} Cite as the starting point. Prop.~4.6 becomes a restatement, not a derivation.

\subsection*{2. Lem.~4.4 --- Partial}

\textbf{Prior.} P.~Tsimpos, N.~Sharp, Y.~Marzouk, \emph{One-Shot Generative Flows: Existence and
Obstructions}, arXiv:2604.15439, Remark~2 with Lemma~24.\Sec

\textbf{Statement reported.} ``Define $\Sigma_t:=\mathrm{Var}(X_t)$ and by Lemma 24 note that we
have $\dot\Sigma_t=\mathrm{Cov}(\dot X_t,X_t)+\mathrm{Cov}(X_t,\dot X_t)$.''

\textbf{Difference.} Their lemma gives the symmetric part of $\mathrm{Cov}(\dot X_t,X_t)$ as
$\tfrac12\dot\Sigma_t$. Lem.~4.4 identifies the antisymmetric complement as $\tfrac12(C^{\T}-C)$,
constant in $t$. For the linear interpolant this follows in two lines from their lemma.

\textbf{Action.} Demote Lem.~4.4 to a remark citing Lemma~24. Do not defend it as a contribution.

\subsection*{3. Prop.~4.9 --- Partial}

\textbf{Prior.} V.~De Bortoli, G.~Deligiannidis, et al., \emph{Augmented Bridge Matching},
arXiv:2311.06978.\F{} C.~Wald, G.~Steidl, \emph{Flow Matching: Markov Kernels, Stochastic
Processes and Transport Plans}, arXiv:2501.16839, Example~4.9.\F

\textbf{Statement quoted.} De Bortoli et al.: ``flow and bridge matching processes preserve the
information of the marginal distributions.'' Wald--Steidl, Ex.~4.9: ``both vector fields generate
the same curves $\mu_t$ but in (d) mass is rotated unnecessarily, which makes the trajectories of
single particles longer than for the vector field in (c).''

\textbf{Difference.} Marginal preservation under arbitrary coupling is stated as background in
both. Prop.~4.9 is its Gaussian linear-map specialisation. What is not in either is the
identification (Rmk.~4.8) of the skew part $N$ as the divergence-free carrier: writing
$\Lambda(t)=V(t)^{-1}NV(t)^{-1}$ one has $\Lambda^{\T}=-\Lambda$ and $NV^{-1}=V\Lambda$, so the
$N$-term preserves every marginal while altering the endpoint map.

\subsection*{4. Thm.~4.11 --- No prior found}

Three independent sources state the general negative --- that different couplings with the same
marginals give different fields --- and none states a criterion for when the endpoint map is
nevertheless invariant.

\textbf{(a) De Bortoli et al., Proposition~2.}\F{} Quoted:
``$\mathbb P^{0,1}_M=\Pi_{0,1}$ if and only if $\Pi_{0,1}=\Pi^{0,1}_*$.''
This is an iff, but about a different object: it characterises when the Markovian projection
\emph{preserves the training coupling}, the answer being that the coupling must already be the
static Schr\"odinger bridge. Theorem~4.11 characterises when the \emph{endpoint map} is
independent of the coupling. Their Gaussian result (Proposition~5,
$\alpha^*=\frac{\sigma^2}{2}(\sqrt{1+4/\sigma^4}-1)$) is the Schr\"odinger-bridge correlation, not
a map-invariance criterion. Retrieved text states that no closed-form characterisation is given
for when different couplings induce identical transport maps.

\textbf{(b) Hertrich et al., Theorem~3(ii).}\F{} Quoted: ``Let $\Sigma_{01}=\Sigma_{10}=0$ and
assume that $\Sigma_0$ and $\Sigma_1$ can be jointly diagonalized.'' The only place they integrate
the field of 3(i) assumes both independence and commuting marginals. No matrix pencil, Frobenius
integrability, zero-curvature argument or time-ordered exponential occurs anywhere in the paper.

\textbf{(c) Tsimpos et al., \S5, open problems.}\Sec{} Quoted: ``Understand intermediate regimes
between independence and determinism. \dots\ It would be very interesting to understand what
becomes possible when the coupling is only partial.'' Theorem~4.11 answers a case of this.

\textbf{Nearest conceptual neighbour.} Wald--Steidl, Corollary~4.8\F{} characterises when
$v_t$ lies in the Wasserstein tangent space, a different property that a reader may conflate with
invariance. The two are independent, and this should be stated in \S4:

For Gaussians a tangent (gradient) field means $K$ symmetric; since $K^{\T}=\tfrac12V^{-1}V'$
under hypothesis (i), $K$ is symmetric iff $[V,V']=0$. Under hypothesis (ii),
$V=p\Sigma_0+q\Sigma_1$ and $V'=p'\Sigma_0+q'\Sigma_1$, so
\[
[V,V']=(pq'-p'q)\,[\Sigma_0,\Sigma_1],
\]
with $pq'-p'q=(2+\sigma+\tau)/4\neq0$ at $t=\tfrac12$ for admissible $(\sigma,\tau)$. Hence under
(i)--(ii) the field is tangent iff $[\Sigma_0,\Sigma_1]=0$. Example~4.18(b) realises invariance
without tangency: endpoint-map drift $6.3\times10^{-5}$ with $\|M-M^{\T}\|=0.165$.

\subsection*{5. Cor.~4.13 --- Same}

\textbf{Prior.} Hertrich et al., Theorem~3(ii).\F{} Under independence and joint
diagonalisability the rectified map is the optimal transport map. Cor.~4.13 is this statement.

\textbf{Action.} Cite. Claim only that Thm.~4.11 shows the hypothesis $[\Sigma_0,\Sigma_1]=0$ is
separate from, and not required for, coupling-invariance (Rmk.~4.12).

\subsection*{6. Cor.~4.21, Rmk.~4.22 --- Partial}

\textbf{Prior.} Tsimpos et al., Proposition~10 with Remark~8.\Sec{} Reported: a field
$v_t(x)=x/(t-\tau)$ with Lipschitz constant $1/|t-\tau|$ diverging as $t\to\tau$, and a process
that ``collapses all points to the origin at time $t=\tau$ and then re-expands them.''

\textbf{Difference.} Their collapse is at an interior time and arises from a rank-deficient
schedule. Cor.~4.21 and Rmk.~4.22 concern collapse at $t=0$ arising from the source law having
strictly smaller support dimension than the target's. The pairing of a non-Lipschitz field with
collapse is theirs; the dimension-counting argument is not stated there.

\subsection*{7. Cor.~5.5 --- Partial}

No specific prior located. Reduction to an invariant $2$-plane by a stabiliser-fixed-point
argument is standard technique in Riemannian flow matching. Present as a lemma, not as a finding.

\subsection*{8. Thm.~5.9 --- No prior found}

No spherical or von Mises--Fisher conditional-flow endpoint law located. Searches run: spherical
flow matching endpoint laws; Riemannian flow matching on $\mathbb S^{d-1}$ with vMF targets;
conditional flow landing point closed form. Negative rests on absence, not on a cleared candidate.

\subsection*{9. Prop.~6.3 --- No prior found}

No prior located on when a conditional flow does or does not improve on returning the source
unchanged. Prop.~6.3 gives the exact criterion $A<\cos j$ and Prop.~6.2 the defining equation.
This and item~12 are the least contested results in the document.

\subsection*{10. \S7, the model $\mathcal M_2$ --- Same construction, opposite conclusion}

\textbf{Prior.} De Bortoli et al., arXiv:2311.06978.\F

\textbf{Statement quoted.} Abstract: they ``augment the velocity field (or drift) with the
information of the initial sample point,'' losing the Markov property but preserving coupling
information. Augmented drift, Eq.~(10):
$d\mathbf X_t=\{b_t(\mathbf X_t)+\sigma_t^2\,\mathbb E[\nabla\log
\mathbb Q_{1|t}(\mathbf X_1|\mathbf X_t)\,|\,\mathbf X_0,\mathbf X_t]\}dt+\sigma_t d\mathbf B_t$,
with sampling initialised at $\mathbf X_0\sim\Pi_0$.

\textbf{Identification.} This is $\mathcal M_2$ of Def.~2.6: the field depends on both the initial
sample and the current state, and sampling starts at the initial sample.

\textbf{Difference.} Their setting is a stochastic bridge with $\sigma_t>0$ and their metric is
coupling recovery; Figure~4 reports improved marginal and coupling scores, and Appendix~E reports
failure only as the pairing becomes increasingly entropic. No fidelity degradation is demonstrated
numerically. The present document studies the deterministic probability-flow ODE with an
informative source and reports the augmentation costing roughly a third of the terminal fidelity
(Meas.~7.6).

\textbf{Action.} Attribute the construction. The contribution is the sign of the effect and its
scaling behaviour, and the paper must be framed against their positive result.

\subsection*{11. Prop.~7.3, Rmk.~7.5 --- Partial}

\textbf{Prior.} Wang et al., arXiv:2412.19992, Theorems~1--2.\Sec{} Reported: ``the limited
performance of pure ODE samplers in diffusion bridge models arises from the singular behavior of
the PF-ODE at the start of the generative process''; Theorem~2, ``At $t=T$, the non-linear drift
term in the PF-ODE (9) becomes singular,'' with the reverse SDE drift remaining well defined. The
mechanism is the Dirac initial condition $q_{T|y}(X_T|y)=\delta(X_T-y)$.

\textbf{Difference.} Prop.~7.3's amplification factor $(1-\alpha t)/(t(1-\alpha))$ and the
non-uniqueness of $\dot y=y/t$ at $y(0)=0$ are sharper statements of a published theorem, not new
ones. Retrieve and read the full text before writing \S7.

\textbf{Adjacent.} Gao et al., arXiv:2605.28962, ``endpoint underfitting'' as $t\to0$;\Sec{}
Zeng \& Yan, arXiv:2509.20952, ``the condition number of the learning problem to diverge'' as
noise levels approach zero, in the noise-source setting.\Sec

\subsection*{12. Meas.~7.6 --- No prior found}

No head-to-head comparison of the two models on a fidelity metric located, and no report of the
inversion. The measured statement is: under maximal-update parameterisation with the learning rate
tuned at width $256$ and transferred, five seeds per cell, and checkpoint selection on a
distributional criterion, the difference is $0.0985$ at width $256$ and $0.1040$ at width $2048$,
while the terminal loss of $\mathcal M_2$ falls from $4.52\times10^{-4}$ to $2.56\times10^{-4}$
and its distributional error from $5.01\times10^{-2}$ to $3.68\times10^{-2}$.

\textbf{Open.} The learning-rate transfer check did not resolve a unique optimum: the
distributional criterion is flat between $10^{-4}$ and $10^{-3}$ at both widths. The sweep should
be repeated at a second learning rate before this is submitted.

\subsection*{13. Cor.~3.3, Rmk.~3.2 --- Same}

\textbf{Prior.} Y.~Blau, T.~Michaeli, \emph{The Perception-Distortion Tradeoff}, CVPR 2018.\Sec

\textbf{Action.} Cite. Claim only the closed-form quantification, Prop.~3.1 and Cor.~3.3, not the
qualitative statement.

\subsection*{14. Stability law refining Thm.~4.11 --- No prior found}

\textbf{Status.} Proved as Proposition~4.28 with Corollary~4.29, measured as Measurement~4.30.
The proposition gives analytic dependence of the endpoint map on the coupling, an explicit
first-order response operator $\mathcal D$, and the fact that $\mathcal D$ annihilates exactly the
directions permitted by hypotheses (i) and (ii); the corollary is the resulting linear law. The
response formula was checked against finite differences of the integrated map: agreement to
$2.4\times10^{-6}$ on a symmetric off-span direction and $5.4\times10^{-6}$ on a skew direction,
with $\|\mathcal D\|$ at the integrator floor $2\times10^{-7}$ on a permitted direction.

\textbf{Design.} Dimensions $n\in\{2,4,8,16,32,64\}$, six independent systems per dimension,
$2048$ random systems per cell, thirty-six coupling configurations per system, all in float64 with
a Heun integrator at $2000$ steps. The clean cell therefore contains $73{,}728$ systems. For each
system a clean coupling is built inside the criterion, then perturbed out of the span of
$\{\Sigma_0,\Sigma_1\}$ by $\varepsilon_{\mathrm{off}}$ and out of symmetry by
$\varepsilon_{\mathrm{skew}}$, each rescaled to keep the joint positive semidefinite.

\textbf{Separation.} The relative displacement of the endpoint map has $90$th percentile
$4.06\times10^{-8}$ across every system satisfying the criterion, against a median of at least
$9.78\times10^{-3}$ once either hypothesis is violated at relative size $0.1$ or above: a
separation of $2.4\times10^{5}$. Marginal matching holds for every coupling tested, violated or
not, at $1.53\times10^{-7}$, confirming Prop.~4.9 carries no hypothesis on the skew part.
Invariance is unaffected by whether the marginals commute, holding at commutator up to $0.199$;
this is the direct refutation of the reading that symmetry of $C$ alone decides the question.

\textbf{The law.} Regressing $\log_{10}$ displacement on $\log_{10}$ violation, $180$ points per
family:
\[
\text{span: slope } 0.9828,\ \text{intercept } -1.0078,\ \text{rms residual } 0.0349;\qquad
\text{skew: slope } 0.9419,\ \text{intercept } -0.1892,\ \text{rms residual } 0.0563.
\]
Both slopes are $1$ to within the fit: the displacement is \emph{linear} in the size of the
violation. Thm.~4.11 is therefore a stability statement, not only a binary criterion, and the two
hypotheses are not interchangeable. At unit relative violation a skew violation displaces the map
by $0.647$ and a span violation by $0.0982$, a ratio of $6.585$.

\textbf{Contradicted prediction.} The register carried the prediction that the two hypotheses are
comparably fragile. It is contradicted and recorded as such. The asymmetry reproduces: an earlier
independent run at smaller scale gave $6.75$.

\textbf{Prior.} No dedicated search was run for a quantitative stability statement of this form.
The criterion it refines has no prior (item~4), so a prior for the refinement is unlikely but not
excluded. Carried in Open.

\subsection*{15. Super-resolution as an instance --- Partial}

\textbf{Status.} New. The identification is exact, not an analogy, and the degenerate-source
hypothesis is measured rather than assumed.

\textbf{The instance.} $c$ is the class label, $z_1$ the high-resolution image, and
$z_0=\mathsf{P}z_1$ with $\mathsf{P}=\mathsf{U}\mathsf{D}$, $\mathsf{D}$ average pooling by $f$ and
$\mathsf{U}$ nearest-neighbour upsampling. $\mathsf{P}$ is an exact projection: measured
null-space energy of the identity arm is $7.5\times10^{-15}$ relative, confirming
$\mathsf{P}^2=\mathsf{P}$ numerically.

\textbf{Why the endpoints stay dependent given $c$.} Conditioned on $c$, $z_0$ is a deterministic
function of $z_1$. The pair is maximally dependent, not independent, which is what Def.~2.1
requires and what an independent-source setting excludes. Conditioning on $c$ removes none of this
dependence, since $c$ is a label and $\mathsf{P}$ acts on the image.

\textbf{Where degeneracy enters.} $\operatorname{rank}\operatorname{Cov}(z_0\mid c)\le
r=3(\mathrm{res}/f)^2$. Measured on STL-10 at $f=4$, job 381593, with strictly more samples than
the bound at every resolution so no estimate is sample-limited:

\begin{center}
\begin{tabular}{@{}rrrrr@{}}
\toprule
res & ambient & bound $r$ & measured rank $z_0$ & samples\\
\midrule
$32$ & $3072$ & $192$ & $192$ & $1024$\\
$48$ & $6912$ & $432$ & $432$ & $1984$\\
$96$ & $27648$ & $1728$ & $1728$ & $3072$\\
\bottomrule
\end{tabular}
\end{center}

Exact equality at all three, against target ranks $1023$, $1983$, $3071$, each sample-limited
full. Hence $\dim\operatorname{supp}(z_0\mid c)<\dim\operatorname{supp}(z_1\mid c)$ strictly, and
Cor.~4.21 applies: a Lipschitz map cannot raise Hausdorff dimension, so a deterministic flow
started at $z_0$ cannot cover the target and collapses to the conditional mean. The measured
null-space energy of the flow confirms the consequence: it restores $0.371$ to $0.433$ of the
target high-frequency energy, and $\mathcal M_1$ lands closer to the least-squares regression
answer than to the target at every resolution.

\textbf{What does not transfer.} The ordering of $\mathcal M_1$ and $\mathcal M_2$ reverses here.
$\mathcal M_2$ scores higher in PSNR at all three resolutions, by $0.114$, $0.062$ and $0.017$,
the last inside one standard error, while restoring $17.5\%$, $12.3\%$ and $11.2\%$ less
high-frequency energy. Between two arms of identical capacity and training that differ only in
whether the network receives $z_0$, higher score with less restored detail is the mean-seeking
inversion of \S~13 rather than a better model. This is reported as a third instance of that
inversion and not as support for the $\mathcal M_1$--$\mathcal M_2$ ordering, which must be
scoped to the sphere.

\textbf{Prior.} M.~Delbracio, P.~Milanfar, \emph{Inversion by Direct Iteration}, TMLR 2023.\Sec{}
Regression to the mean in deterministic restoration is their stated motivation, and Blau \&
Michaeli give the perception-distortion form. \textbf{Shared:} the phenomenon.
\textbf{Not shared:} the identification of super-resolution as an exact instance of Def.~2.1 with
the rank bound stated and measured, and the dimension-theoretic obstruction as the mechanism
rather than an MMSE argument.

\textbf{Action.} Claim the instance and the measured rank equality. Do not claim the phenomenon.

\subsection*{16. Prop.~4.33 --- Partial}

\textbf{Status.} Proved as Proposition~4.33 with Definition~4.31 and Lemma~4.32, measured as
Measurement~4.34, read in Remark~4.35.

\textbf{Statement.} The value of the flow-matching objective at its population optimum is not a
function of the endpoint map. The witness is explicit and one-dimensional: for $n=1$ and
$\Sigma_0=\Sigma_1=1$ every cross-covariance $c\in[-1,1]$ satisfies both hypotheses of
Thm.~4.11 and gives $M=1$, while the loss floor runs from $\mathcal F(1)=0$ to
$\mathcal F(0)=\pi/2$. Hence no modulus of continuity in $\|M(C_1)-M(C_2)\|$ can bound
$|\mathcal F(C_1)-\mathcal F(C_2)|$.

\textbf{Prior.} Y.~Zhang et al., \emph{Bridge Graphical Models}, arXiv:2608.19144, Appendix~K.\F{}
\textbf{Shared:} the object. Their Markovization gap
$\mathfrak G(\mathsf Q)=\int_0^1\mathbb E[\operatorname{Tr}\operatorname{Var}(U_t\mid X_t)]dt$
is Definition~4.31, and
they report the coupling moving it by two orders of magnitude. \textbf{Not shared:} that this
quantity is decoupled from the endpoint map, which requires Thm.~4.11 and has no prior. The
measurement sharpens it: on couplings that are provably equivalent at the endpoint the floor
spans a median factor $2.23$, while on couplings that move the map by a median $63\%$ it spans
only $1.23$. The floor is ordered against the map, not merely uninformative about it.

\textbf{Action.} Cite Zhang et al. for the object and for the two-orders-of-magnitude report.
Claim the decoupling and the ordering. This row also discharges the obligation recorded under
that entry in \S~Checked and cleared, that their result be distinguished explicitly so a reader
does not take it as a contradiction.

\section{Checked and cleared}

\textbf{Cheng \& Schwing, C2OT, ICCV 2025, arXiv:2503.10636.}\F{} The source is a Gaussian prior
throughout; the training objective is
$\mathbb E_{t,(x_0,x_1,c)\sim\pi}\|v_{\theta,c}(t,x_t)-u(x_t|x_0,x_1)\|^2$, so the network never
receives $x_0$. Their formal content (Eqs.~7--9) is that minibatch optimal transport skews
$q(x_0|c)\neq q(x_0)$ at training while test-time sampling is unbiased. No analysis of
ill-conditioning, determinacy, or degeneracy at $t=0$. Different setting and different mechanism.

\textbf{Tsimpos et al., arXiv:2604.15439, main result.}\Sec{} The object is straightness of the
flow under endpoint independence, characterised as requiring a deterministic coupling. Not
endpoint-map invariance. Cleared as the main result; \emph{not} cleared on Rmk.~2, Lem.~24,
Thm.~9, Rmk.~5, Prop.~10, Rmk.~8, which appear above.

\textbf{Zhang et al., \emph{Bridge Graphical Models}, arXiv:2608.19144.}\F{} Verified to exist via
the arXiv API, dated 2026-06-16; full text retrieved. The object is the Markovization gap
$\mathfrak G(\mathsf Q)=\int_0^1\mathbb E[\mathrm{Tr}\,\mathrm{Var}(U_t|X_t)]dt$, the irreducible
floor of the flow-matching loss, not the endpoint map. Their Appendix~K reports the coupling
changing that gap by two orders of magnitude; this is consistent with Thm.~4.11 and must be
distinguished explicitly in the text, since a reader may take it as a contradiction. Their
Theorem~2, $W_2(\rho_1^\theta,\rho_1)\le e^L\epsilon$, is a candidate scaffold for a bound in
\S7, subject to regularity assumptions that fail in the present setting.

\textbf{Lim, arXiv:2512.16768.}\F{} Studies flow matching through finite-sample plug-in
estimation. Proposition~5.5 parameterises flux-null corrections by antisymmetric matrices
$A_i(t)^{\T}=-A_i(t)$ and requires each $\Sigma_i(t)$ symmetric positive definite, excluding
degenerate sources. Equivalence is by probability-flux divergence, not by coupling. Relevant to
Rmk.~4.8 as the general form of the $N$-term; not a prior for Thm.~4.11.

\textbf{Wald \& Steidl, arXiv:2501.16839.}\F{} A survey. Example~2.7 gives the Gaussian optimal
transport map; Proposition~4.11 treats the independent coupling with standard Gaussian source. No
general-coupling covariance computation of the endpoint map, no matrix pencil, no conditioning of
the velocity network on the initial sample. Example~4.10 treats Dirac data and mass-splitting at a
crossing time, not lower-dimensional source support.

\textbf{Pooladian et al., Multisample Flow Matching, PMLR v202.}\Sec{} Coupling design to
straighten flows while preserving marginals. Different object.

\section{Open}

\begin{enumerate}\itemsep2pt
\item Wang et al., arXiv:2412.19992: retrieve the full text and read Theorems~1--2 before writing
\S7. Item~11 rests on a secondary summary.
\item Tsimpos et al.: retrieve Lemma~24, Remark~2, Theorem~9, Remark~5, Proposition~10 and
Remark~8 directly. Items~2 and 6 rest on secondary summaries.
\item Tong et al., arXiv:2302.00482, and Tong et al., \emph{Simulation-Free Schr\"odinger Bridges
via Score and Flow Matching}, PMLR v238: not checked. The general statement that any coupling with
correct marginals yields a valid probability path is likely to be located in the former, which
bears on item~3.
\item Blau \& Michaeli: confirm the citation directly.
\item Guidance behaviour and the discrete-diffusion threshold: not searched. Report both as
observations without a priority claim, or search before claiming novelty.
\item Image-to-image bridge literature (I2SB, DDBM, InDI, Direct Diffusion Bridges,
stochastic interpolants, Doob $h$-transform): not swept individually. This is where an empirical
report of the item~10 effect is most likely to exist.
\item Item~12: repeat the width sweep at a second learning rate.
\item Item~14: the numbered statement is written, Prop.~4.28 and Cor.~4.29. Still owed: one
targeted search for quantitative perturbation bounds on Gaussian rectified-flow endpoint maps
before the phrase ``no prior found'' is used in the submission.
\item Item~16: the $47$-system measurement is at $n\in\{3,6\}$ with the loss floor evaluated in
closed form rather than trained. The dimension sweep of Meas.~4.30 would extend it to $n=64$ at
small cost, since the same random-system machinery applies.
\item Item~15: retrieve Delbracio \& Milanfar in full and determine whether a dimension or
Lipschitz obstruction is stated there, or only the conditional-mean argument. The verdict Partial
depends on this.
\item Item~15: closed. The rank equality is now exact at $\mathrm{res}=32$, $48$ and $96$, none
sample-limited.
\item Item~15: the null-space energy ratio is not a collapse index on its own. The regression arm
records the \emph{highest} ratio, $0.68$ to $0.78$, together with the lowest PSNR of the learned
arms, so the statistic also counts incorrect high-frequency energy. The comparison is sound only
between arms of matched capacity and training, which is why it is drawn for
$\mathcal M_1$ against $\mathcal M_2$ and for nothing else. A distributional metric on the
restoration outputs would settle the ordering; none was run.
\end{enumerate}

\end{document}
